\documentclass[11pt]{article}

\usepackage[margin=2.2cm]{geometry}
\usepackage{amsmath}
\usepackage{amssymb}
\usepackage{enumitem}
\usepackage{titlesec}
\usepackage{xcolor}
\usepackage{fancyhdr}

\titleformat{\section}{\large\bfseries}{\thesection}{0.6em}{}
\titlespacing{\section}{0pt}{1.4em}{0.8em}

\pagestyle{fancy}
\fancyhf{}
\fancyhead[L]{Solving Polynomial Equations}
\fancyhead[R]{Name: \makebox[4cm]{\hrulefill}}
\fancyfoot[C]{\thepage}
\renewcommand{\headrulewidth}{0.4pt}

\newcommand{\worksheettitle}{Solving Polynomial Equations}
\newcommand{\qspace}{\vspace{1.1cm}}

\setlist[enumerate,1]{leftmargin=1.6em, itemsep=0.4em, topsep=0.2em}

\begin{document}

\begin{center}
  {\Large\bfseries \worksheettitle}\\[0.3em]
  {\normalsize Fluency \(\cdot\) Problem Solving \(\cdot\) Reasoning \(\cdot\) Enrichment}
\end{center}

\vspace{0.5em}

\section{Fluency}

\subsection*{Solving equations already in factored form}
For each equation, use the zero-factor law to find all values of $x$.

\begin{enumerate}
  \item $(x-2)(x+3) = 0$
  \qspace
  \item $(2x+1)(3x+7) = 0$
  \qspace
\end{enumerate}

\subsection*{Factorising monic quadratics}
Factorise each expression, then solve the resulting equation.

\begin{enumerate}[resume]
  \item $x^2 - 7x + 12 = 0$
  \qspace
  \item $x^2 + 3x - 10 = 0$
  \qspace
  \item $x^2 - 4x - 21 = 0$
  \qspace
\end{enumerate}

\subsection*{Factorising non-monic quadratics}
Factorise each expression by splitting the middle term, then solve the resulting equation.

\begin{enumerate}[resume]
  \item $2x^2 + 7x + 3 = 0$
  \qspace
  \item $3x^2 - 5x - 2 = 0$
  \qspace
  \item $6x^2 + 11x - 10 = 0$
  \qspace
\end{enumerate}

\subsection*{Rearranging before factorising}
Rearrange each equation so that one side is equal to zero, then factorise and solve.

\begin{enumerate}[resume]
  \item $3x^2 + 2x = 5$
  \qspace
  \item $2x^2 - 3x = 9$
  \qspace
\end{enumerate}

\newpage
\section{Problem Solving}

\begin{enumerate}
  \item The volume of water in a bucket after $t$ minutes is modelled by
  \[
    V = t^2 + 5t + 6,
  \]
  where $V$ is measured in litres. How many minutes does it take for the bucket to contain 20 litres of water? (Discard any solution that does not make sense.)
  \qspace

  \item A rectangular garden bed has a length that is 3 m longer than its width. The area of the garden bed is 40 m$^2$. Let $w$ be the width, in metres. Write an equation for the area and solve it to find the width of the garden bed.
  \qspace

  \item A ball is thrown upward from ground level. Its height above the ground, $h$ metres, after $t$ seconds is given by
  \[
    h = -5t^2 + 20t.
  \]
  Find the times at which the ball is 15 m above the ground. Explain why there are two possible answers.
  \qspace

  \item A company's weekly profit, $P$ (in thousands of dollars), from selling $x$ hundred items is modelled by
  \[
    P = x^2 - 2x - 15.
  \]
  Find the number of items (to the nearest hundred) that must be sold for the company to break even, i.e.\ $P = 0$.
  \qspace

  \item The product of two consecutive positive integers is 132. Let $n$ be the smaller integer. Write an equation and solve it to find both integers.
  \qspace

  \item A picture frame is designed so that its length is 1 cm more than twice its width. The area of the frame is 55 cm$^2$. Let $w$ be the width, in centimetres. Write an equation for the area and solve it to find the width of the frame.
  \qspace
\end{enumerate}

\newpage
\section{Reasoning}

\begin{enumerate}
  \item State the \textbf{zero-factor law} (also called the zero product property) in your own words.
  \qspace

  \item Explain \emph{why} the zero-factor law works. In your explanation, consider what it means for two numbers to multiply to give zero, and why at least one of them must be zero.
  \qspace

  \item Consider the equation $(x-2)(x+3) = 6$. Explain why we \emph{cannot} apply the zero-factor law directly to this equation, even though the left-hand side is already factored. What must be done to the equation before it can be solved by factorising?
  \qspace

  \item When factorising a non-monic quadratic such as $2x^2 + 7x + 3$, we ``split the middle term'' into two terms. Explain why this method works, and why the two numbers we look for must multiply to give $a \times c$ and add to give $b$.
  \qspace
\end{enumerate}

\newpage
\section{Enrichment}

Some equations involving powers can be solved by rewriting them as polynomial equations. If every term can be written using the \emph{same base}, the equation can be rewritten as a quadratic in a new variable (for example, $y = 3^x$), which can then be solved by factorising.

\begin{enumerate}
  \item Solve for $x$:
  \[
    9^x - 4(3^x) + 3 = 0.
  \]
  \emph{Hint: write $9^x$ as a power of 3, then let $y = 3^x$.}
  \qspace

  \item Solve for $x$:
  \[
    4^x - 6(2^x) + 8 = 0.
  \]
  \emph{Hint: write $4^x$ as a power of 2, then let $y = 2^x$.}
  \qspace
\end{enumerate}

\end{document}
